% ===================================================================
%  TABEL Nr. 2 — Die menslike liggaam. — Die pouse. — Die speletjies.
%  Meme regime que le tableau 1 : traduction de 2026, cles de
%  texto/io, renvois aux memes objets, gras sur le terme seul.
%
%  LA NEGATION SE FERME, ET C'EST LA REGLE QUE L'OUTIL NE VOIT PAS.
%  « nie ... nie » : la seconde negation clot la proposition, et une
%  phrase neerlandaise a negation simple reste correcte a l'oeil de
%  kolonoj.py. C'est le point ou il faut relire — ce tableau en compte
%  onze.
%
%  ET LE PLURIEL AFRIKAANS N'EST PAS CELUI DU NEERLANDAIS : « oë »,
%  « ore », « lippe », « ribbes », « longe » ; le tremas y marque
%  l'hiatus (oë, knieë) et non un emprunt.
%
%  LES NOMS DE JEUX SONT NATIONAUX, la note (*) du fac-simile le dit
%  elle-meme : l'ido les a forges. On prend le nom afrikaans quand le
%  jeu existe — blindemol, kennetjie, bok-bok —, et une description
%  quand il n'existe pas.
% ===================================================================

%%K t02-apar-1 apar
\VUsaut{4.0mm}
\VUtitre{15.0pt}{60}{TABEL Nr. 2}
\VUsaut{2.0mm}
\VUtitre{11.4pt}{0}{Die menslike liggaam. --- Die pouse.}
\VUtitre{11.4pt}{0}{Die speletjies.}

%%K t02-tit-0 sub
\VUtitre{13.2pt}{0}{Die menslike liggaam.}
\VUsaut{2.4mm}
\VUcentre{10.2pt}{0}{\textit{(Eenvoudige les in die natuurwetenskap.)}}
\VUsaut{3.0mm}

%%K t02-tit-1 sub pos=t02-01-1
\VUpk{10.2pt}{40}{Die kop.}
\VUsaut{3.0mm}

%%K t02-01-1 p
1. --- Die vernaamste dele van die menslike liggaam is: die \VUgras{kop}, die
\VUgras{romp} en die \VUgras{ledemate}.

%%K t02-02-1 p
2. --- Die kop bevat twee dele: die \VUgras{skedel}, wat die \VUgras{brein} bevat en
wat die \VUgras{hare} bedek, en die \VUgras{gesig}.

%%K t02-03-1 p
3. --- Op ons tekening sien ons nie die skedel of die brein nie; maar ons sien baie
duidelik die belangrike dele van die gesig.

%%K t02-04-1 p
4. --- Hier is eers die \VUgras{voorkop}, wat 'n kragtige verstand aandui, 'n denke wat
altyd werk. Die \VUgras{slape} is baie vry. Die \VUgras{oog} is lewendig en wyd oop. 'n
Digte \VUgras{wenkbrou} en lang \VUgras{ooghare} beskerm dit.

%%K t02-05-1 p
5. --- Hier is verder die \VUgras{neus} en die \VUgras{neusgate}. Die neus dien ons om
te ruik en om asem te haal. Aan weerskante van die neus is die \VUgras{wange}, verder
weg die \VUgras{ore}. Die onderste deel van die gesig bestaan uit die \VUgras{mond} en
die \VUgras{ken}.

%%K t02-06-1 p
6. --- Ons sien hulle nie baie goed nie, want die \VUgras{snor} en die
\VUgras{wangbaard} verberg hulle vir ons. Maar ons weet dat die mond die twee
\VUgras{lippe} het, die \VUgras{bokaak} en die \VUgras{onderkaak} met hulle
\VUgras{tande}, en die \VUgras{tong}. Met die mond praat en eet ons. Met die tong proe
ons die voedsel.

%%K t02-07-1 p
7. --- Die oog, die oor, die neus en die mond is, net soos die vingers, die organe van
die vyf sintuie: die \VUgras{gesig}, die \VUgras{gehoor}, die \VUgras{reuk}, die
\VUgras{smaak}, die \VUgras{tas}.

%%K t02-08-1 p
8. --- Die kop is dus een van die interessantste dele van die menslike liggaam.

%%K t02-tit-2 sub
\VUsaut{4.4mm}
\VUpk{10.2pt}{40}{Die romp.}
\VUsaut{3.0mm}

%%K t02-09-1 p
9. --- Die \VUgras{nek} verbind die kop met die romp. Dit bevat die \VUgras{keel} en
die \VUgras{agternek}.

%%K t02-10-1 p
10. --- Die romp bevat ook baie noodsaaklike organe. In die \VUgras{bors} is die
\VUgras{longe}, waarmee ons asem haal, en die \VUgras{hart}, wat die middelpunt van die
lewe is. Wanneer die hart ophou klop, sterf die lewende wese.

%%K t02-11-1 p
11. --- Die \VUgras{ribbes} hou die bors op.

%%K t02-12-1 p
12. --- Die ander dele van die romp is die \VUgras{sy}, die \VUgras{rug}, die
\VUgras{skouers} en die \VUgras{buik}. Ons lê op ons sy of op ons rug om te slaap, ons
dra die vragte op ons skouers.

%%K t02-13-1 p
13. --- In die buik is die \VUgras{maag} en die \VUgras{ingewande}. Hulle dien ons om
die voedsel te verteer.

%%K t02-tit-3 sub
\VUsaut{4.4mm}
\VUpk{10.2pt}{40}{Die ledemate.}
\VUsaut{3.0mm}

%%K t02-14-1 p
14. --- Ons het vier \VUgras{ledemate}: twee \VUgras{arms} en twee \VUgras{bene}. Met
die arms neem, hou, lig en versamel ons die voorwerpe; met die bene staan, loop en hardloop
ons.

%%K t02-15-1 p
15. --- Die \VUgras{skouer} verbind die arm met die liggaam. 'n Baie sterk
\VUgras{spier}, die biseps, beweeg die arm wat die \VUgras{elmboog} met die
\VUgras{voorarm} verbind. Die \VUgras{gewrig} vorm die lit van die \VUgras{hand}, wat
eindig met die \VUgras{vingers} en die \VUgras{naels}. Op die hand onderskei ons 'n
\VUgras{aar} (en 'n slagaar) waarin die \VUgras{bloed} vloei.

%%K t02-16-1 p
16. --- Die dele van die been is: die \VUgras{dy}, die \VUgras{knie} en die
\VUgras{knieholte}, die \VUgras{kuit}, waarvan die \VUgras{vleis} met die \VUgras{vel}
bedek is, die \VUgras{enkel} en die \VUgras{voet}. Die \VUgras{bene} van die been is
baie dik en baie sterk. Aan die punt van die voet is die \VUgras{tone}. Die ander dele
is die \VUgras{voetsool} en die \VUgras{hakskeen}.

%%K t02-17-1 p
17. --- So is die byna volledige beskrywing van die menslike liggaam.

%%K t02-17-2 p
\textit{Nota.} --- \textit{Met behulp van die uitdrukking: « ek het pyn aan... (die kop
ens.) » sal die onderwyser hierdie hele woordeskat weer kan gebruik uit die oogpunt van
die goeie of slegte} \VUgras{liggaamstoestand}.

%%K t02-tit-4 sub
\VUsaut{5.0mm}
\VUpk{10.2pt}{40}{Gimnastiek.}
\VUsaut{2.4mm}
\VUcentre{10.2pt}{0}{\textit{(Vertelling deur een van die leerlinge.)}}
\VUsaut{3.0mm}

%%K t02-18-1 p
18. --- Om lenig, rats en gesond te wees, moet ons ons bene en ons arms oefen. Daarom
speel ons en beoefen ons \VUgras{gimnastiek}.

%%K t02-19-1 p
19. --- Gedurende die week vind hierdie twee oefeninge op die binnehof van die liseum
plaas (sien tabel nr. 1). Maar Donderdae en Sondae gaan ons na 'n groot \VUgras{grasveld},
waar ons plek het om te hardloop en ons te vermaak.

%%K t02-20-1 p
20. --- Ons meesters en ons ouers vergesel ons soms daarheen; maar dit is die
\VUgras{gimnastiekonderwyser}\textsuperscript{(12)} wat die oefeninge lei.

%%K t02-21-1 p
21. --- Op die grasveld, links van die ingang, staan die \VUgras{gimnastiektoestelle}
\textsuperscript{(1)}, ons klim teen die \VUgras{leer}\textsuperscript{(2)} op, ons
swaai onderstebo aan die \VUgras{ringe}\textsuperscript{(3)} en aan die
\VUgras{sweefstok}\textsuperscript{(4)}, ons trek ons met die hande op tot bo aan die
\VUgras{touleer}\textsuperscript{(5)} of aan die \VUgras{tou met
knope}\textsuperscript{(6)}.

%%K t02-22-1 p
22. --- Intussen spring ons maats oor die \VUgras{houtperd}\textsuperscript{(7)} of oor
die \VUgras{brug met gelyke leggers}\textsuperscript{(11)}. Ander
\VUgras{boks}\textsuperscript{(8)} of \VUgras{skerm}\textsuperscript{(9)} met maskers en
met 'n \VUgras{floret}. Laastens \VUgras{lig ander gewigte}\textsuperscript{(10)}.

%%K t02-tit-5 sub
\VUsaut{4.6mm}
\VUpk{10.2pt}{40}{Die speletjies.}
\VUsaut{3.0mm}

%%K t02-23-1 p
23. --- Rondom die turners speel seuntjies en dogtertjies allerlei \VUgras{speletjies}.
Die dogtertjies vermaak hulle met hulle \VUgras{hoepel}\textsuperscript{(14)}; hulle
spring \VUgras{tou}\textsuperscript{(13)} of nooi hulle broers en neefs vir 'n partytjie
\VUgras{kroukie}\textsuperscript{(15)}. Hulle is verruk wanneer hulle die \VUgras{bal}
van die teenstander wegstoot op die oomblik dat hy dink hy sal maklik deur die hekkie
gaan!

%%K t02-24-1 p
24. --- Die seuns verkies die speletjies wat meer krag vra. Hulle speel graag
\VUgras{kegels}\textsuperscript{(16)} wat hulle behendig met 'n \VUgras{bal}
\textsuperscript{(17)} omgooi. Hulle versmaai ook nie die \VUgras{tol}
\textsuperscript{(18)} en die \VUgras{balletjie-en-bekertjie}\textsuperscript{(19)} of
selfs die \VUgras{padda}\textsuperscript{(20)} nie. Maar hulle gunstelingspeletjies is
die \VUgras{albasters}\textsuperscript{(21)} en die \VUgras{muurbal}
\textsuperscript{(22)}, want die muur van die \VUgras{broeikas}\textsuperscript{(26)} is
baie geskik om die ligte bal te laat terugspring.

%%K t02-25-1 p
25. --- Die klein Jan, wat op sy \VUgras{stelte}\textsuperscript{(24)} loop, is baie
behendig en weet baie goed hoe om sy ewewig te hou. Naby hom speel 'n paar maats
\VUgras{wegkruipertjie}\textsuperscript{(25)} in daardie broeikas en agter die
\VUgras{heining}. Ander verkies die \VUgras{raai-wie-geslaan-het}\textsuperscript{(28)}
of die \VUgras{pluimbal}\textsuperscript{(29)} wat van die een \VUgras{raket} na die
ander vlieg.

%%K t02-noto-1 noto
\VUnotes{143.2mm}{%
(*) Die kinder- en seunsspeletjies het dikwels suiwer nasionale name. Ons het probeer om
hulle so goed as moontlik in Ido te noem, óf deur ons deur die handeling self te laat
lei, wat hulle kenmerk, óf deur die nasionale naam uit hierdie of daardie land te kies,
wanneer dit nie te onbillik is nie. Bv.: die \VUgras{vatspel} (in Duits en in Frans) is
die \VUgras{paddaspel} \textsuperscript{(20)} (in Engels), waarin die spelers probeer om
hulle ronde skyfies deur die bek van die metaalpadda te gooi wat met oop bek op die
deurboorde kis (vat) staan. Die \VUgras{skouersprong}\textsuperscript{(30)} verskil van
die \VUgras{rugsprong} net hierin: om oor die kop te spring, steun die spelers hulle
hande op die skouers van hulle maat, terwyl hulle in die « \VUgras{rugsprong} » hulle
hande net op sy rug steun. --- Of ook is 'n paar van die deelnemers gebuk terwyl die
ander oor hulle rug sal spring met die hande op die skouer; die eerstes moet die skok
verdra sonder om in te gee, die tweedes moet spring sonder om hulle ewewig te verloor
(sien die tabel self).%
}

%%K t02-26-1 p
26. --- In die middel van die grasveld staan twee bome. Aan die voet van een van die
twee het sewe of agt seuns 'n partytjie \VUgras{skouersprong}\textsuperscript{(30)} (*)
gereël. Hierdie spel is nie in alle lande bekend nie; maar almal weet wat die
\VUgras{rugsprong} is, wat die seuns hierdie keer nie speel nie.

%%K t02-tit-6 sub
\VUsaut{5.0mm}
\VUpk{10.2pt}{40}{Die speletjies in die buitelug.}
\VUsaut{3.4mm}

%%K t02-27-1 p
27. \VUgras{Blindemol}\textsuperscript{(31)} is ook baie vermaaklik; dogtertjies en
seuntjies sit om die beurt die \VUgras{blinddoek} oor hulle oë en gryp die onhandiges
wat hulle laat vang.

%%K t02-28-1 p
28. --- In die middel van die grasveld, hier is 'n baie Franse maar effens ruwe spel:
dit is die \VUgras{beer}\textsuperscript{(32)}, wat veel meer lawaai maak as die so
rustige spel van die \VUgras{vlieër}\textsuperscript{(33)} of selfs as die Engelse spel
\VUgras{tennis}\textsuperscript{(34)}.

%%K t02-29-1 p
29. --- Hierdie laaste sport is die vreugde van die groot leerlinge. Ons onderwysers
self beoefen dit graag. Dit vra groot behendigheid en ratsheid en dit beval my baie.
Ek laat die spel van die \VUgras{swaai}\textsuperscript{(35)} aan die dogtertjies oor.

%%K t02-30-1 p
30. --- Ek hou nie van 'n ander Engelse spel wat maklik gevaarlik sou word nie.

%%K t02-30-1b p
Ek bedoel die \VUgras{voetbal}\textsuperscript{(36)} wat my ouer broer bly maak. Ek
verkies ons ou \VUgras{aanjaagspel}\textsuperscript{(38)}, net so gesond en werklik
minder ru.

%%K t02-tit-7 sub
\VUsaut{4.6mm}
\VUpk{10.2pt}{40}{Fietsry en balonspel.}
\VUsaut{3.0mm}

%%K t02-31-1 p
31. --- Ek ry ook graag op my \VUgras{fiets}\textsuperscript{(37)} om die grasveld.

%%K t02-32-1 p
32. --- Volgende jaar sal ek \VUgras{perdry}\textsuperscript{(39)} leer. Hoe mooi is 'n
perd wat sierlik stap of draf, en wat in galop oor die hindernisse spring!

%%K t02-33-1 p
33. --- In die somer leer ons \VUgras{swem}\textsuperscript{(40)}. Ons duik en bad met
genot in die helder en vars water van die rivier. Soms laat ons op die ou end 'n papier
\VUgras{ballon}\textsuperscript{(41)} op, wat statig in die lug styg sodra dit met warm
lug gevul is.

%%K t02-34-1 p
34. --- Daar is ook nog baie ander speletjies, maar ek sou nie klaarkom om hulle op te
noem nie, en uit hierdie opsomming sien 'n mens dat ons op Donderdae nie te veel verveeld
raak op ons mooi grasveld nie.

%%K t02-34-2 p
L.W. --- \textit{Die onderwyser sal hierdie hoofstuk maklik aanvul deur elkeen van die
belangrikste speletjies uitvoeriger te beskryf.}
